\subsection*{Drill: Spending Points}
\label{drill:points/spending}

This drill requires enough blockers to challenge a jammer, and someone emulating a bench coach.

This drill focuses on blockers looking to `spend' their points to hold back the jammer.

Use either a two, three or four wall to match the jammer to a point where they are unlikely to escape without offence.
A three or four wall is preferred.

The jammer is free to jam and has between 10 and 40 seconds to score as many points as they can.
The jammer may not know how long they have remaining, all they know is that when their bench says to call they have to call. 


This emulates either clearing a skater from the penalty box, or trying to score points before the other jammer reaches the wall.  
The jammer is on a timer and does not expect to complete a pass beyond this one.

The blockers are looking not to protect their points, but to spend them efficiently. 
The idea being that if a jammer has the full fourty seconds to score points while the wall tries to hold the score at zero, the jammer is more likely to take the whole wall. 
The most efficient way to prevent the jammer from scoring points is to take them off the track and recycle them. 


To achieve this the wall should recognise whose point has already been taken, and hence who is now free to take an inferior position and drive the jammer.

\begin{itemize}
\item If all points are still in play, then the wall should either look to spend one point for a drive, or hold the jammer scoreless. 
\item If the jammer has scored a point, that blocker is now `free' to drive the jammer without concern for position.
\item If the jammer has scored two points then more elaborate four wall pockets are possible.
\end{itemize}


{\it Progression: } Add an offence to support the jammer. 
